\chapter{CÁC GIẢI PHÁP VÀ ĐÓNG GÓP NỔI BẬT}

Chương này trình bày ba giải pháp kỹ thuật nổi bật được phát triển trong quá trình xây dựng hệ thống OMR, tập trung vào ba thách thức cụ thể của ảnh chụp điện thoại trong môi trường giáo dục thực tế. Phần đầu mô tả cơ chế tự động hiệu chỉnh lưới câu hỏi MCQ Map Search Rescue, giải quyết vấn đề lệch lưới sau khi warp phối cảnh. Phần hai trình bày thuật toán tính điểm bong bóng tổng hợp kết hợp \texttt{density} và \texttt{darkness}, cải thiện độ nhạy với bong bóng tô mờ. Phần cuối mô tả giải pháp Smart Camera Scanner trên trình duyệt, cho phép chụp và chấm phiếu trực tiếp từ camera điện thoại không cần cài ứng dụng. Hai giải pháp đầu được đánh giá định lượng bằng ablation study trên bộ dữ liệu 61 ảnh.

\section{Cơ chế tự động hiệu chỉnh lưới câu hỏi -- MCQ Map Search Rescue}
\label{sec:mcq-rescue}
\subsection{Vấn đề gặp phải}

Trong thực tế, giáo viên chụp ảnh phiếu trong nhiều điều kiện khác nhau: ánh đèn huỳnh quang chênh lệch, ảnh chụp hơi nghiêng, học sinh tô không đủ đậm hoặc có vết tẩy xóa. Sau khi warp phiếu về ảnh chuẩn, lưới MCQ đôi khi vẫn bị lệch vài pixel theo chiều dọc hoặc khoảng cách dòng, làm nhiều câu bị đọc là không chắc chắn (\path{uncertain_count} cao) hoặc tô nhiều đáp án giả.

\subsection{Nguyên nhân kỹ thuật}

Sau bước warp phối cảnh, hệ thống tính khoảng cách dòng trung bình bằng cách phát hiện hai hàng marker đặc trưng ở đầu và cuối vùng MCQ:
$$\texttt{line\_h} = \frac{y_{\text{cuối}} - y_{\text{đầu}}}{\text{số khoảng}}$$

Công thức này giả định tất cả hàng câu cách đều nhau. Giả định đó chỉ đúng khi phép warp bù hoàn hảo mọi biến dạng của tờ giấy. Thực tế tồn tại hai nguồn sai số có tính hệ thống:

\begin{itemize}[leftmargin=1.5cm]
  \item \textbf{Warp phối cảnh 4 điểm không bù được độ cong nhẹ của giấy.}
    Biến đổi \texttt{getPerspectiveTransform} ánh xạ 4 góc phát hiện lên hình chữ nhật chuẩn, nhưng giấy thực tế luôn cong nhẹ khi đặt bàn hoặc cầm tay. Sau warp, khoảng cách dọc giữa các hàng không hoàn toàn đều mà biến thiên nhỏ theo chiều cao.
  \item \textbf{Sai số phát hiện góc tích lũy theo số thứ tự câu.}
    Mỗi góc marker phát hiện với sai số khoảng $\pm 1$\,pixel; sai số này truyền vào ma trận warp và ảnh hưởng đến \texttt{line\_h}. Hàm decode dùng ngoại suy tuyến tính \texttt{cy = top\_center\_y + row\_idx $\times$ line\_h} — câu 1 hầu như không bị ảnh hưởng, nhưng câu 15–20 có thể lệch 8–12\,pixel so với tâm bong bóng thực tế.
\end{itemize}

Hàm \path{_decode_mcq_with_map()} lấy mẫu cửa sổ $\pm 0{,}44 \times \texttt{line\_h}$ quanh \texttt{cy}. Khi \texttt{cy} lệch quá 5–6\,pixel, cửa sổ trượt ra khỏi tâm bong bóng và điểm giảm xuống dưới ngưỡng nhận diện, khiến câu bị đánh dấu không chắc chắn -- dù học sinh đã tô rõ ràng.

\subsection{Cách giải quyết}

Cơ chế Rescue được triển khai trong \path{be/app/services/omr/omr_service.py} tại hàm \path{process_omr_exam()}, chạy ngay sau lần decode MCQ baseline.

\subsubsection*{Điều kiện kích hoạt}

Cơ chế Rescue chỉ được kích hoạt khi số câu không chắc chắn trong kết quả baseline vượt ngưỡng:
$$\text{gate} = \max\!\left(4,\; \operatorname{round}(0.40 \times Q)\right)$$
trong đó $Q$ là tổng số câu hỏi. Với đề 40 câu, gate $= 16$; với đề 80 câu, gate $= 32$. Ngưỡng 40\% số câu đảm bảo Rescue chỉ can thiệp khi có dấu hiệu lệch lưới đáng kể, không kích hoạt với phiếu chỉ có vài câu tô không rõ. Ngoài ra, cơ chế bị tắt tự động ở chế độ \textit{long-form} (phiếu $\geq 25$ câu/block hoặc $\geq 4$ block) và có thể tắt thủ công qua \texttt{disable\_mcq\_rescue} trong profile.

\subsubsection*{Thuật toán tìm kiếm lưới cục bộ}

Khi điều kiện kích hoạt thỏa mãn, hệ thống thực hiện tìm kiếm toàn diện trên không gian hai chiều của tham số hình học lưới:

\begin{lstlisting}[caption={Không gian tìm kiếm MCQ Map Search Rescue}, label={lst:rescue}, basicstyle=\ttfamily\small, breaklines=true, frame=single, showstringspaces=false]
line_scales    = [1.00, 0.92, 0.88, 1.08, 0.84, 1.16]   # 6 ty le line-height
shift_mults    = [0.0, -0.5, 0.5, -1.0, 1.0, -1.5, 1.5, -2.0, 2.0]  # 9 top-shift

for scale in line_scales:                       # 6 gia tri
    cand_line_h = line_h * scale               # dieu chinh khoang cach dong
    for shift_mul in shift_mults:              # 9 gia tri
        cand_shift = shift_mul * cand_line_h   # dich vi tri dong dau tien
        cand = decode_mcq(line_h=cand_line_h, top_shift=cand_shift)
        cand_rank = (uncertain(cand), double_mark(cand), -quality(cand))
        if cand_rank < best_rank:
            best_rank, best_result = cand_rank, cand
# Tong: 6 x 9 = 54 ung vien
\end{lstlisting}

Tổng không gian tìm kiếm là $6 \times 9 = 54$ ứng viên. Mỗi ứng viên gọi lại đầy đủ hàm \texttt{\_decode\_mcq\_with\_map()}, tức là thực hiện toàn bộ quá trình tính điểm bong bóng và phân loại kết quả cho tất cả câu hỏi.

\subsubsection*{Hàm xếp hạng ứng viên}

Mỗi ứng viên được xếp hạng theo bộ ba từ điển (tuple comparison theo Python):
$$\text{rank} = \bigl(\text{uncertain\_count},\ \text{double\_mark\_count},\ -\text{quality}\bigr)$$

Trong đó hàm chất lượng được tính:
$$\text{quality}(\text{kết quả}) = 10 \times |\text{câu chắc chắn}| + 6 \times |\text{4 câu đầu chắc chắn}| + 70 \times \sum \text{margin}_q$$

Bộ ba này ưu tiên giảm số câu không chắc chắn trước, sau đó giảm tô nhiều đáp án, cuối cùng tối đa điểm chất lượng tổng. Hệ số 70 trên \texttt{margin\_sum} làm cho kết quả có margin cao (tô rõ ràng) được ưu tiên hơn kết quả nhiều câu vừa đạt ngưỡng tối thiểu.

\subsubsection*{Tiêu chí chấp nhận bảo thủ}

Ứng viên tốt nhất chỉ thay thế baseline khi đáp ứng một trong ba điều kiện sau:
\begin{itemize}[leftmargin=1.5cm]
  \item Giảm ít nhất 2 câu không chắc chắn ($\Delta\text{uncertain} \geq 2$); hoặc
  \item Giảm ít nhất 1 câu không chắc chắn và không tăng số câu tô nhiều đáp án; hoặc
  \item Số câu không chắc chắn không đổi nhưng giảm được số câu tô nhiều đáp án.
\end{itemize}

Nếu không có ứng viên nào đạt tiêu chí, kết quả baseline gốc được giữ nguyên. Tiêu chí chấp nhận bảo thủ này ngăn hệ thống thực hiện điều chỉnh không cần thiết: nếu phiếu thực sự tô không rõ hoặc tô sai (không phải lỗi lưới), không có cấu hình lưới nào có thể cải thiện đáng kể, và baseline được bảo toàn.

\subsubsection*{Hiệu quả và chi phí}

Khi Rescue thành công, hệ thống ghi mã cảnh báo \texttt{MCQ\_MAP\_SEARCH\_RESCUE} vào \texttt{warning\_codes} và lưu metadata đầy đủ: \texttt{line\_h\_before}/\texttt{after}, \texttt{top\_shift\_px}, \texttt{initial\_uncertain}, \texttt{final\_uncertain}, \texttt{tested} (số ứng viên đã thử). Điều này cho phép audit chính xác cơ chế đã điều chỉnh gì.

Về chi phí tính toán, trong trường hợp xấu nhất Rescue thử 54 lần decode MCQ trên cùng một phiếu. Tuy nhiên cơ chế này chỉ kích hoạt khi phiếu đã vượt ngưỡng uncertain, tức là phiếu tốt (uncertain $< 60\%$ số câu) không phát sinh chi phí nào. Trên thực tế, phần lớn phiếu chụp đúng cách đều có uncertain thấp và bỏ qua hoàn toàn bước này.

\begin{figure}[H]
  \centering
  \includegraphics[width=\textwidth,height=0.72\textheight,keepaspectratio]{figures/MCQ_Map_Search_Rescue.jpg}
  \caption{Luồng MCQ Map Search Rescue}
  \label{fig:mcq-map-search-rescue}
\end{figure}

\subsection{Kết quả sau xử lý}

Khi ảnh bị lệch nhẹ và số câu không chắc chắn vượt ngưỡng, hệ thống có thêm cơ hội sửa lưới MCQ thay vì trả ngay kết quả baseline. Nếu rescue cải thiện được rank, số câu không chắc chắn giảm và bảng đáp án trả về ổn định hơn; nếu không cải thiện, baseline vẫn được giữ lại để tránh thay đổi sai. Test liên quan là \texttt{TC-GRADE-02}, trong đó ảnh có bong bóng tô mờ phải ghi \path{uncertain_count} và kích hoạt rescue khi vượt ngưỡng.

Để đo tác động cụ thể, hệ thống được eval trên 61 ảnh với hai chế độ: có và không có Rescue. Kết quả ở Bảng~\ref{tab:rescue-ablation} cho thấy Rescue cải thiện G-DetRate hơn 26 điểm phần trăm và giảm số câu uncertain trung bình từ 5{,}43 xuống còn 0{,}13 mỗi phiếu.

\begin{longtable}{|L{3.2cm}|L{2.5cm}|L{2.5cm}|L{2.2cm}|L{2.2cm}|}
\caption{Ablation: ảnh hưởng của MCQ Map Search Rescue đến độ chính xác nhận diện (61 ảnh, 4 nhóm)}
\label{tab:rescue-ablation}\\
\hline
\textbf{Cấu hình} & \textbf{G-DetRate} & \textbf{Unc /phiếu} & \textbf{Shadow} & \textbf{Tilted} \\
\hline
\endfirsthead
\hline
\textbf{Cấu hình} & \textbf{G-DetRate} & \textbf{Unc /phiếu} & \textbf{Shadow} & \textbf{Tilted} \\
\hline
\endhead
Không có Rescue & 72{,}9\% & 5{,}43 & 70{,}7\% & 66{,}0\% \\
\hline
Có Rescue & \textbf{99{,}3\%} & \textbf{0{,}13} & \textbf{99{,}5\%} & \textbf{100{,}0\%} \\
\hline
\textbf{Cải thiện} & $+26{,}4\,\text{pp}$ & $-5{,}30$ & $+28{,}8\,\text{pp}$ & $+34{,}0\,\text{pp}$ \\
\hline
\end{longtable}

\subsection{Ví dụ minh họa -- Trace Rescue thực tế}

Ví dụ dưới đây mô tả luồng xử lý đầy đủ khi một phiếu 40 câu kích hoạt Rescue, dựa trên cấu trúc metadata mà hệ thống ghi vào \path{bubble\_confidence\_json}.

\textbf{Đầu vào:} Phiếu thi 40 câu, chụp bằng điện thoại dưới đèn huỳnh quang lệch về phía phải. Sau bước warp phối cảnh, lưới MCQ bị lệch xuống khoảng 8--10 pixel theo chiều dọc so với tọa độ lý tưởng trong profile.

\textbf{Decode baseline:}
\begin{itemize}[leftmargin=1.5cm]
  \item \texttt{uncertain\_count = 26} $>$ gate $= \max(4,\ \lfloor 0.40 \times 40 \rceil) = 16$.
  \item Điều kiện kích hoạt Rescue thỏa mãn; hàm \path{process\_omr\_exam()} bắt đầu tìm kiếm cục bộ.
\end{itemize}

\textbf{Quá trình Rescue:}
\begin{itemize}[leftmargin=1.5cm]
  \item Thử 54 ứng viên ($6 \times 9$); mỗi ứng viên gọi đầy đủ \path{\_decode\_mcq\_with\_map()}.
  \item Ứng viên tốt nhất tìm được: \texttt{scale = 0.92} (rút ngắn khoảng cách dòng 8\%), \texttt{shift\_mul = -0.5} (dịch lưới lên trên $0.5 \times \texttt{line\_h}$).
  \item Rank ứng viên: $(4,\ 0,\ -\text{quality})$ so với baseline $(26,\ 2,\ -\text{quality})$ -- ứng viên tốt hơn rõ ràng.
\end{itemize}

\textbf{Kết quả:}
\begin{itemize}[leftmargin=1.5cm]
  \item \texttt{uncertain\_count} sau Rescue $= 4$ (giảm 22 câu so với baseline).
  \item \texttt{double\_mark\_count = 0}.
  \item \texttt{warning\_codes}: \texttt{["MCQ\_MAP\_SEARCH\_RESCUE"]}.
  \item Metadata ghi: \texttt{line\_h\_before=18.5}, \texttt{line\_h\_after=17.0}, \texttt{top\_shift\_px=-9.2}, \texttt{tested=54}.
\end{itemize}

Nếu không có cơ chế Rescue, 26 câu sẽ bị trả về dạng không chắc chắn và điểm số sẽ không phản ánh đúng bài làm của học sinh.

\section{Thuật toán tính điểm bong bóng tổng hợp}
\subsection{Vấn đề gặp phải}

Phương pháp đơn giản nhất để xác định bong bóng đã tô là đếm số pixel đen trong ô bằng hàm \texttt{countNonZero}. Tuy nhiên, phương pháp này có độ chính xác thấp khi bong bóng tô mờ hoặc khi nền ảnh không hoàn toàn trắng sau khi nhị phân hóa.

\subsection{Nguyên nhân kỹ thuật}

Ảnh camera có ánh sáng không đều, bóng giấy, nhiễu từ đường kẻ và nét tô không đồng nhất. Threshold nhị phân có thể làm mất nét tô mờ hoặc phóng đại vùng nhiễu. Nếu chỉ dựa vào số pixel đen, một ô bị đường kẻ chạy qua hoặc nền tối cục bộ có thể có score gần giống ô thật sự được tô.

\subsection{Cách giải quyết}

Trong \path{be/app/services/omr/omr_mcq.py}, hàm \texttt{\_cell\_score()} kết hợp tỷ lệ pixel nhị phân (\texttt{density}) và độ tối từ ảnh xám gốc (\texttt{darkness}) vào một điểm tổng hợp. Công thức tính và các đại lượng thành phần đã được mô tả đầy đủ tại Mục~\ref{sec:pipeline-buoc-13}. Ngoài điểm ô đơn lẻ, hệ thống còn dùng hàm \texttt{\_row\_candidate\_quality()} để chọn đáp án tốt nhất trong mỗi hàng câu hỏi:

\begin{equation}
\begin{aligned}
\mathrm{quality} ={}& 1.10 \times \mathrm{best\_score}
+ 1.45 \times \mathrm{margin} \\
&+ 0.25 \times \mathrm{best\_center}
+ 0.12 \times \mathrm{best\_dark}
- \mathrm{left\_penalty}
\end{aligned}
\end{equation}

Trong đó \texttt{margin} là khoảng cách giữa điểm cao nhất và nhì, \texttt{left\_penalty} phạt nếu ô đầu tiên, tương ứng với cột A, có điểm cao nhưng biên độ quá nhỏ. Trường hợp này thường là nhiễu từ đường kẻ trái.

\begin{longtable}{|L{3.4cm}|L{10.6cm}|}
\caption{Ý nghĩa các đại lượng trong công thức chấm bong bóng}
\label{tab:bubble-score-fields}\\
\hline
\textbf{Đại lượng} & \textbf{Ý nghĩa} \\
\hline
\texttt{density} & Tỷ lệ pixel đen trong vùng bong bóng sau nhị phân hóa. \\
\hline
\texttt{darkness} & Độ tối tổng hợp từ ảnh xám, giúp giữ tín hiệu tô mờ. \\
\hline
\texttt{margin} & Khoảng cách giữa đáp án tốt nhất và đáp án đứng thứ hai. \\
\hline
\texttt{best\_score} & Điểm density của lựa chọn tốt nhất. \\
\hline
\texttt{best\_center} & Mật độ vùng trung tâm, giảm nhiễu từ viền ô. \\
\hline
\texttt{best\_dark} & Độ tối bổ sung từ ảnh xám. \\
\hline
\texttt{left\_penalty} & Hình phạt cho nhiễu ở cột A khi biên độ không đủ rõ. \\
\hline
\end{longtable}

\subsection{Kết quả sau xử lý}

Cách tính mới giúp hệ thống ổn định hơn với bong bóng tô nhẹ, vì ảnh xám gốc vẫn đóng góp một phần vào score khi ảnh nhị phân bị mất nét. Đồng thời \texttt{margin}, \texttt{best\_center} và \texttt{left\_penalty} giúp hạn chế chọn nhầm do đường kẻ hoặc vết bẩn. Test liên quan gồm \texttt{TC-GRADE-01} với ảnh phiếu chuẩn và \texttt{TC-GRADE-02} với ảnh có câu không rõ.

Một bổ sung quan trọng trong module chấm điểm là điều kiện \textit{strong\_signal\_bypass}: nếu \texttt{best\_score} $\geq 0{,}55$ và \texttt{best\_center} $\geq 0{,}72$, kết quả được chấp nhận ngay cả khi noise gate của hàng câu đó bình thường sẽ đánh dấu uncertain. Logic này ngăn chặn trường hợp bong bóng tô rõ, tâm bong bóng đặc, bị từ chối nhầm do noise gate quá chặt với một số câu hỏi cụ thể (Q7, Q9, Q11). Bảng~\ref{tab:noise-gate-ablation} so sánh kết quả trên 58 ảnh đánh giá trước và sau khi thêm strong\_signal\_bypass.

Để đánh giá đóng góp của thành phần \texttt{darkness} trong công thức điểm, hệ thống được thử nghiệm bổ sung với chế độ \textit{density-only} (bỏ qua \texttt{darkness}, chỉ dùng \texttt{density}). Bảng~\ref{tab:scoring-ablation} cho thấy cả hai công thức đạt cùng độ chính xác nhận diện câu, nhưng weighted scoring giảm được thời gian xử lý trung bình từ 9{,}7 xuống 5{,}4 giây mỗi phiếu - do công thức tổng hợp giúp phân biệt bong bóng tô nhẹ rõ hơn, hạn chế kích hoạt Rescue không cần thiết.

\begin{longtable}{|L{3.5cm}|L{2.5cm}|L{2.5cm}|L{2.5cm}|L{2.5cm}|}
\caption{Ablation: ảnh hưởng của \textit{strong\_signal\_bypass} đến độ chính xác nhận diện câu (58 ảnh)}
\label{tab:noise-gate-ablation}\\
\hline
\textbf{Cấu hình} & \textbf{G-DetRate} & \textbf{Uncertain /phiếu} & \textbf{Tô rõ} & \textbf{Shadow} \\
\hline
\endfirsthead
\hline
\textbf{Cấu hình} & \textbf{G-DetRate} & \textbf{Uncertain /phiếu} & \textbf{Tô rõ} & \textbf{Shadow} \\
\hline
\endhead
Không có bypass & 89{,}9\% & 2{,}12 & 79{,}1\%\newline(4{,}79 unc) & 95{,}5\%\newline(0{,}90 unc) \\
\hline
Có bypass & \textbf{97{,}8\%} & \textbf{0{,}43} & \textbf{94{,}2\%}\newline(1{,}16 unc) & \textbf{99{,}5\%}\newline(0{,}10 unc) \\
\hline
\textbf{Cải thiện} & $+7{,}9\,\text{pp}$ & $-1{,}69$ & $+15{,}1\,\text{pp}$ & $+4{,}0\,\text{pp}$ \\
\hline
\end{longtable}

\begin{longtable}{|L{4.5cm}|L{2.5cm}|L{2.2cm}|L{4.5cm}|}
\caption{Ablation: density-only vs weighted scoring trên 61 ảnh (kết hợp với Rescue)}
\label{tab:scoring-ablation}\\
\hline
\textbf{Công thức score} & \textbf{G-DetRate} & \textbf{Unc/phiếu} & \textbf{Thời gian TB} \\
\hline
\endfirsthead
\hline
\textbf{Công thức score} & \textbf{G-DetRate} & \textbf{Unc/phiếu} & \textbf{Thời gian TB} \\
\hline
\endhead
\texttt{density} (countNonZero) & 99{,}3\% & 0{,}13 & 9{,}7\,s/ảnh \\
\hline
\texttt{0,90×density + 0,10×darkness} (hiện tại) & \textbf{99{,}3\%} & \textbf{0{,}13} & \textbf{5{,}4\,s/ảnh} \\
\hline
\end{longtable}

\subsection{Ví dụ minh họa -- So sánh countNonZero và công thức tổng hợp}

Ví dụ dưới đây minh họa sự khác biệt giữa hai phương pháp trên cùng một câu hỏi, dựa trên cấu trúc dữ liệu confidence JSON mà hệ thống ghi lại.

\textbf{Tình huống:} Câu 15 trên phiếu có học sinh tô lựa chọn B nhưng lực tô nhẹ (bút bi màu nhạt, đường kẻ in chạy qua ô B). Sau bước nhị phân hóa, một phần nét tô bị mất do ngưỡng Otsu tính trên toàn ảnh hơi cao.

\textbf{Các giá trị đo được tại ô B, câu 15:}
\begin{itemize}[leftmargin=1.5cm]
  \item \texttt{density} (tỷ lệ pixel đen sau nhị phân hóa): $0.11$
  \item \texttt{darkness} (độ tối tổng hợp từ ảnh xám gốc):
    $0.55 \times 0.14 + 0.30 \times 0.25 + 0.15 \times 0.31 = 0.20$
  \item \texttt{score} tổng hợp: $0.90 \times 0.11 + 0.10 \times 0.20 = 0.119$
\end{itemize}

\textbf{Kết quả so sánh:}

\begin{longtable}{|L{5.5cm}|L{2.5cm}|L{5.8cm}|}
\caption{So sánh hai phương pháp tại câu 15}
\label{tab:vi-du-bubble-score}\\
\hline
\textbf{Phương pháp} & \textbf{Quyết định} & \textbf{Lý do} \\
\hline
\texttt{countNonZero} đơn thuần & \textbf{Blank} & density $= 0.11 <$ ngưỡng $0.15$ \\
\hline
Công thức tổng hợp & \textbf{B} (đúng) & score $= 0.119 >$ dynamic\_mark; darkness từ ảnh xám xác nhận vùng trung tâm tối hơn nền \\
\hline
\end{longtable}

Ví dụ này cho thấy khi ảnh nhị phân mất một phần tín hiệu tô mờ, thành phần \texttt{darkness} từ ảnh xám gốc đóng vai trò cứu vãn, giúp hệ thống không bỏ qua đáp án hợp lệ. Ngược lại, nếu đường kẻ in gây ra \texttt{density} giả ở cột A, \texttt{left\_penalty} và \texttt{margin} nhỏ sẽ chặn lựa chọn sai này.

\section{Smart Camera Scanner trên trình duyệt}

Để chụp và chấm phiếu nhanh chóng từng tờ một, giáo viên cần giao diện camera trực tiếp trên điện thoại không cần cài ứng dụng gốc. Việc phát hiện khi nào phiếu đã được đặt đúng vị trí, không nghiêng và đủ bốn góc là thách thức kỹ thuật cần giải quyết hoàn toàn phía client.

Trình duyệt chỉ cho phép dùng camera qua \texttt{navigator.mediaDevices.getUserMedia}~\cite{webrtc} trong secure context như HTTPS hoặc localhost. Frontend không thể chạy toàn bộ pipeline OpenCV nặng như backend để chấm ngay trên video frame. Vì vậy cần một bộ kiểm tra nhẹ ở client để quyết định khi nào ảnh đủ điều kiện chụp: nhìn thấy giấy, đủ bốn marker tối, marker có tương phản với nền và trạng thái ổn định qua nhiều frame.

Hệ thống triển khai thuật toán phát hiện marker tại hàm \texttt{evaluateAlignment()} trong file \path{MultichoicePage.tsx}. Mỗi 130ms, thuật toán phân tích một frame video:

\begin{enumerate}
    \item Vẽ frame lên canvas ẩn và đọc pixel data.
    \item Tính \texttt{centerLuma}, là độ sáng trung bình vùng giữa viewfinder để xác nhận đang nhìn thấy tờ giấy trắng.
    \item Tính \texttt{darkLumaThreshold = clamp(centerLuma - 26, 28, 110)}, là ngưỡng thích nghi theo điều kiện ánh sáng.
    \item Tại bốn vị trí góc được cấu hình trong Form Profile, tính \texttt{darkRatio}.
\end{enumerate}

\begin{equation}
\texttt{lum\_pixel}
=
0.299 \times R
+
0.587 \times G
+
0.114 \times B
\end{equation}

\begin{equation}
\texttt{darkRatio}
=
\frac{\texttt{count}(\texttt{lum} < \texttt{darkLumaThreshold})}
{\texttt{totalPixels}}
\end{equation}

Trạng thái \texttt{LOCKED} chỉ được kích hoạt khi đồng thời thỏa mãn ba điều kiện:

\begin{equation}
\texttt{hasFourDarkMarkers}
=
\texttt{all}(\texttt{darkRatio}[c] \geq 0.14)
\end{equation}

\begin{equation}
\texttt{paperInsideFrame}
=
\texttt{centerLuma} \geq 52
\end{equation}

\begin{equation}
\texttt{markerContrastOk}
=
(\texttt{centerLuma} - \texttt{mean(markerLuma)}) \geq 20
\end{equation}

Nếu bốn frame liên tiếp đều thỏa mãn các điều kiện trên, hệ thống chuyển sang trạng thái \texttt{locked} và khung hướng dẫn trên giao diện chuyển sang màu xanh. Nếu trạng thái \texttt{locked} được duy trì trong 1300ms, hệ thống tự động chụp ảnh dưới định dạng JPEG với chất lượng \texttt{quality = 0.92}, sau đó gửi ảnh đến API để chấm điểm.

Luồng camera dùng \texttt{navigator.mediaDevices.getUserMedia} với camera sau nếu có. Frame video được vẽ lên canvas ẩn, sau đó frontend chỉ tính toán các đại lượng đơn giản như \texttt{centerLuma}, \texttt{darkRatio} và \texttt{markerLuma}. Cách làm này đủ nhẹ để chạy trên điện thoại và vẫn tận dụng cùng Form Profile với backend, vì vị trí marker được lấy từ profile đang chọn.

Smart Camera Scanner cho phép giáo viên căn phiếu trong khung và để hệ thống tự chụp khi đủ điều kiện. Khi môi trường không phải HTTPS/localhost, frontend hiển thị lỗi rõ ràng và người dùng vẫn có thể dùng chế độ tải ảnh từ thư viện. Test liên quan là \texttt{TC-CAM-01}, kiểm tra video hiển thị, trạng thái \texttt{searching}/\texttt{locked} và điều kiện tự chụp.
